\cotefile{157-2}
\scanpage{179}
\mspage{\unsure{124}}
% archive sheet number: 178

% slanted
\margin{\uline{NB} Si on suppose $f$ propre on doit pouvoir prendre aussi $\unsure{\mathcal F|F}=\mathrm{Cat}_{/S}$.}

\deleted{\unsure{Pour|Donc} \unsure{supposer} \unsure{qu'il} \unsure{soit} \unsure{lisse} \illegible{1} \unsure{$X\mapsto X\times_S S'$} \unsure{transforme} \unsure{$W_S^d$} \unsure{en} $\unsure{W_{S'}^{d*}|W_{S'}^d}$.}
b) Le foncteur \unsure{chgt.|cht.} de base \deleted{\unsure{défini par} \unsure{$f:S'\to S$}, \unsure{dans}}

\[
\underline{\mathrm{Liss}}\deleted{\unsure{{}_{/S}}}\,S\longrightarrow
\underline{\mathrm{Liss}}\deleted{\unsure{{}_{/S'}}}\,S',\qquad
X\longmapsto X\times_S S'\mathbin{/}S'
\]

transforme $W_S$ en $W_{S'}$.

b$'$) Énoncé analogue pour $\underline{\mathrm{Fib}\,S}\longrightarrow\underline{\mathrm{Fib}}\,S'$.

b$''$) Énoncé analogue pour $\underline{\mathrm{Fibsc}}\,S\longrightarrow\underline{\mathrm{Fibsc}}\,S'$.

\uline{Dém} Si $\mathcal F$ est une des trois catégories $\deleted{\mathrm{Cat}_{/S}}\inserted{\underline{\mathrm{Liss}}\,S}$, $\underline{\mathrm{Fib}}\,S$, ou $\underline{\mathrm{Fibsc}}\,S$, le foncteur $\rho_S$ peut se décrire comme le foncteur de localisation canonique

\[
\deleted{\unsure{\mathrm{Fib}}}\ \mathcal F(W_S^d)^{-1}\longrightarrow\mathcal F W_S^{-1}
\]

(ça marcherait pour n'importe quelle sous-catégorie pleine $\mathcal F$ de $\mathrm{Cat}_{/S}$ qui contient $\underline{\mathrm{Fib}}\,S$). Or \deleted{\unsure{si on}} \inserted{si b), b$'$) ou b$''$)} est satisfait, on trouve un diagramme commutatif de foncteurs induits par le changement de base

\diagram{\begin{tikzcd}[ampersand replacement=\&,column sep=large]
\mathcal F(W_S^d)^{-1} \arrow[r,"\rho_S"] \arrow[d,"{f^*_{\deleted{\unsure{\mathcal F}}W_S^d}}"'] \& \mathcal F(W_S)^{-1} \arrow[d,"{f^*_{W_S}}"] \\
\mathcal F'(W_{S'}^d)^{-1} \arrow[r,"\rho_{S'}"'] \& \mathcal F'(W_{S'})^{-1}
\end{tikzcd},}

\deleted{On sait} sous réserve, il est vrai, dans le cas $\mathcal F=\mathrm{Cat}_{/S}$, que les foncteurs image inverse soient aussi compatibles avec $W_S^d$, $W_{S'}^d$ (ce qui est automatique si $\mathcal F=\underline{\mathrm{Fib}}\,S$ ou $\underline{\mathrm{Fibsc}}\,S$). On sait que $f^*_{W_S^d}$ décrit le foncteur $f'$ sur les types d'homotopie relative, et \deleted{\unsure{peut}} prouve a) comme conséquence
\scanpage{180}
\mspage{125}
% archive sheet number: 179

de b \unsure{(toutefois)|(resp.)}, b$'$, ou b$''$\deleted{\unsure{b|)}}, il reste à s'assurer que le foncteur $f^*_{W_S}$ décrit de $=$ le foncteur induit par le précédent sur les types d'hom. relatifs loc. \unsure{constants}. Or on a les inclusions canoniques $\unsure{\mathcal F|F}_0\subset\mathcal F$, $\mathcal F'_0\subset\mathcal F'$, et sur $\mathcal F_0,\mathcal F'_0$ $W_S=W_S^d\underset{\unsure{2}}=W_S^{\unsure{t|f}}$ et ces \deleted{\unsure{quatre}} catégories sont OK pour \deleted{calculer} décrire $\mathrm{HOT}^{\unsure{\mathrm{lc}|\mathrm{le}}}_{W}(S\deleted{\unsure{\,) }})$ et $\mathrm{HOT}^{\mathrm{lc}}_W(S')$, et le foncteur $\unsure{g^*|f^*}$ induit entre eux par images inverses.

% slanted
\margin{NB $(\mathrm{Cat}_{/S})_0=\underline{\unsure{\mathrm{Fibhot}|\mathrm{Fiblct}}}\,S$, pour les autres cas c'est clair.}

(\uline{NB} Le foncteur chgt. de base $X\mapsto X\times_S S'$ est compatible avec les localiseurs sur ces catégories, car c'est trivial sous la forme $W_S^f$.)

On a alors, par images inverses \inserted{$X\mapsto X\times_S S'$}, un diagramme commutatif

\diagram{\begin{tikzcd}[ampersand replacement=\&,column sep=large]
\mathcal F_0(W_S)^{-1}\arrow[r,"\approx"]\arrow[d,"{(f^*_{W_S})_0}"'] \& \mathcal F_{\deleted{\unsure{0}}}(W_S)^{-1}\arrow[d,"{f^*_{W_S}}"]\\
\mathcal F'_0(W_{S'})^{-1}\arrow[r,"\approx","{\deleted{\unsure{\simeq}}}"'] \& \mathcal F'(W_{S'})^{-1}
\end{tikzcd},}

où les flèches horizontales sont des \uline{équivalences}\unsure{\,/s} de catégories, par le \unsure{th. satellite} 1 \textcircled{I}. Cela prouve que $f^*_{W_S}$ lui aussi est bon pour calculer les images inverses de t. d'hom. relatifs.

\deleted{\unsure{On}} On a prouvé que \deleted{b$''$ \unsure{$\Longrightarrow$} b} b, b$'$, b$''$ impliquent chacun a), prouvons que a les implique. Or a) dit exactement que les images inverses \deleted{\unsure{par}} $g^*$ des éléments de $\mathrm{HOT}^{\mathrm{lc}}_W(S)$
\scanpage{181}
\mspage{126}
% archive sheet number: 180

sont, à isom. près, \deleted{\unsure{les}} obtenus par le foncteur image inverse sur $\unsure{\mathcal F|F}$ \deleted{\unsure{tout}} entier, en interprétant $\mathrm{HOT}^{\unsure{\mathrm{lc}|\mathrm{le}}}_W(S)$ comme $\mathcal F W_S^{-1}$. Mais cela implique bien que \deleted{les flèches de $\mathcal F$ qui \unsure{sont des}} \inserted{$\mathcal F\longrightarrow\mathcal F'$ doit être} compatible avec les localiseurs. Mais je \unsure{devrais|devine} -- on trouve seulement que si $u:X\longrightarrow Y$ de $\mathcal F$ est dans $W_S$, alors la flèche $u':X'\longrightarrow Y'$ donne un iso dans $\mathcal F'(W_{S'})^{-1}$. \deleted{\unsure{Je ne}} Je \unsure{ne} sais pas, d'autre part, si cela implique que $u'$ est dans $W_{S'}$.

\uline{Corollaire 1} : Supposons que $p:S'\longrightarrow S$ soit une Hot-fibration. Alors le foncteur image inverse $X\mapsto X\times_S S'$ de $\mathrm{Cat}_{/S}$ dans $\mathrm{Cat}_{/S'}$ applique $W_S$ dans $W_{S'}$, \deleted{$W_S^d$ dans $\unsure{W_{S'}^d|W_{S'}^t}$ (et $W_S^{\unsure{g|q}}$ dans $W_{S'}^g$)}, donc le foncteur $p^*:\mathrm{HOT}_W(S)\longrightarrow\mathrm{HOT}_W(S')$ commute aux foncteurs de localisation $\rho_S,\rho_{S'}$ dans ces catégories. (\uline{NB} Il faut faire l'hypothèse adéquate sur $W$, savoir L 6 bis page 44. C'est OK pour $W_\infty$.)

Le fait que $X\mapsto X\times_S S'$ soit compatible avec $W_S,W_{S'}$ est une tautologie, compte tenu de l'hypothèse L 6 bis. Il \unsure{ne|me} semble pas qu'il \deleted{\unsure{convienne}} \inserted{soit} aussi compatible à $W_S^d,W_{S'}^d$ (je devrais fabriquer un contreexemple, \deleted{\unsure{avec}} bien sûr $S'$ non propre sur $S$), \deleted{donc b) \unsure{ne me} sert pas directement. Mais} \inserted{Donc} on est sous les conditions de \unsure{b)|b$'$)}, qui implique a).
\scanpage{182}
\mspage{127}
% archive sheet number: 181

Je ne me sens pas capable d'exhiber un foncteur autre qu'une Hot-fibration (ou \deleted{\unsure{tout}} au moins une $W$-fibration), pour lequel le foncteur image inverse sur $\mathrm{HOT}_W$ soit compatible aux foncteurs de localisation $\rho_S,\rho_{S'}$. Bien sûr, si $X$ est quelconque dans $\mathrm{Cat}_{/S}$, rien n'empêche de regarder son localisé pour $W_S$, i.e. $\rho_S(\mathrm{hot}_{W,S}(X))$, et de \deleted{\unsure{se|de}} demander comment trouver son image inverse. Si $p:S'\longrightarrow S$ est une Hot-fibration, alors il suffit de prendre $\mathrm{hot}_{W,S'}(X')$, $X'=X\times_S S'$. Dans \deleted{tout} autre cas, (j'entends sans hyp. sur $X/S$, \deleted{\unsure{sans supposer}} plus précisément sans supposer que $X$ est lui-même Hot-fibré sur $S$), cette formule est \unsure{essentiell\textsuperscript{t}|considérée} sans \unsure{espoir|emploi|objet}. On doit, pour calculer $f^*\rho_S(\mathrm{hot}_{W,S}(X))$, ``calculer'' $\rho_S(\mathrm{hot}_{W,S}(X))$ en ``montant à l'étage'', comme $\mathrm{hot}_{W,S}(C(X))$. On trouve alors l'image inverse $\mathrm{hot}_{W,S'}(C(X)\times_S S')$ (v\textsuperscript{o} \unsure{cor.} 1 p 119).

\centerline{\(\backslash\)}

Maintenant que la loi contravariante de $\mathrm{HOT}_W(S)$ en $S\in\mathrm{Ob}\,\mathrm{Cat}$ semble bien comprise, il s'agit de voir qu'on a \inserted{bien} un dérivateur $\mathrm{HOT}$. Je ne m'attends pas à ce que les $\mathrm{HOT}^{\unsure{\mathrm{lc}|\mathrm{le}}}_W(S)$
\scanpage{183}
\mspage{128}
% archive sheet number: 182

définissent eux aussi un dérivateur, mais il faudra m'en assurer expressément. (C'est pour l'existence de $f_!,f_*$ que ça doit clocher.)

À dire vrai, \deleted{\unsure{il}} il \deleted{\unsure{fau}}\unsure{faudrait} pour bien faire regarder aussi la structure de \uline{2-foncteur} de

\[
S\longmapsto\mathrm{HOT}_W(S)\qquad \mathrm{Cat}^{\unsure{\mathrm o|\circ}}\longrightarrow\mathrm{Cat}
\]

donc pour $S,S'$ données dans $\mathrm{Cat}$, définir

\[
\underline{\mathrm{Hom}}(S',S)^{\mathrm o}\longrightarrow
\underline{\mathrm{Hom}}(\mathrm{HOT}_W(S),\mathrm{HOT}_W(S')).
\]

Cel\unsure{a} revient \unsure{seulement}, pour

\[
 f,g:S'\rightrightarrows S
\]

deux flèches de $S'$ dans $S$, et un hom de foncteurs

\[
u:f\longrightarrow g,
\]

à décrire un hom de foncteurs

\[
u^*:g^*\longrightarrow f^*
\]

\deleted{\unsure{de}}\unsure{entre|d'entre} $\mathrm{HOT}_W(S)$ et $\mathrm{HOT}_W(S')$. C'est d'ailleurs chose faite, dans le cadre général \unsure{de|des} $(\mathcal M,W)$, quand on décrit $\mathrm{HOT}_{W}(S)$ et $\mathrm{HOT}_W(S')$ par localisation de $\underline{\mathrm{Hom}}(S^{\mathrm o},\mathrm{Cat})$ et $\underline{\mathrm{Hom}}(S'^{\mathrm o},\mathrm{Cat})$. Ces catégories s'interprètent resp. comme $\underline{\mathrm{Fibsc}}\,S\deleted{\unsure{s}}$ et $\underline{\mathrm{Fibsc}}\,S'$, donc comme des catégories de catégories sur $S$ resp. $S'$, on voudrait donc interpréter $u^*$ en termes de telles catégories, et tant qu'à faire, pour des catégories $X$ sur $S$ aussi générales que possible. \deleted{\unsure{Quand on}} \inserted{Nous nous}
\scanpage{184}
\mspage{129}
% archive sheet number: 183

attendons donc à devoir faire des hypothèses, soit sur $X$, soit sur $f$ et sur $g$, qui assurent que les images de $\mathrm{hot}_{W,S}(X)$ par $f^*$ et $g^*$ ne sont rien d'autres que $\mathrm{hot}_{W,S}(f^*_{\deleted{\unsure{S}}}X)$, $\mathrm{hot}_{W,S}(g^*X)$ respectivement, où on pose

\[
f^*(X)=X\times_S(S',f),\quad g^*(X)=X\times_S(S',g).
\]

Donc on \deleted{supposera} devra sans doute se placer dans l'une des hypothèses suivantes :

a) $X$ est lisse sur $S$

b) $X$ est Hot-fibrée sur $S$

c) $f$ et $g$ sont propres (et sous réserve de la validité de la conjecture page 122).

On s'attend alors à ce que $u^*$ soit associé\deleted{\unsure{e}} à un hom

\[
u^*:g^*(X)\longrightarrow f^*(X)\qquad\text{\uline{dans}}\ \mathrm{Cat}_{/S'}
\]

défini en termes de $u$. Y \deleted{\unsure{a}} a-t-il moyen de le définir bel et bien, sans supposer $X$ \deleted{\unsure{f}} au moins fibrée sur $S$? Prenons $S'=e$, donc $f,g$ correspondent à des objets $s,t\in\mathrm{Ob}\,S$, et $u$ à une flèche

\[
u:s\longrightarrow t,
\]

les images inverses de $X/S$ par $f,g$ sont les fibres $X_s,X_t$, et on cherche $\deleted{\unsure{X_t}}$

\[
u^*:X_t\longrightarrow X_s.
\]

On n'a aucune chance d'en trouver \inserted{canoniquement,} si on ne suppose $X$ fibrée sur $S$! Donc ce ne sont pas a), b) ou c) qui peuvent \unsure{nous} servir, mais unique
\scanpage{185}
\mspage{130}
% archive sheet number: 184

ment l'hypothèse $X\in\mathrm{Ob}\,\underline{\mathrm{Fib}}\,S$. Donc on ne peut espérer \unsure{donner|trouver} une description de $u^*$ \inserted{\deleted{\unsure{sur}} $\mathrm{HOT}_W(S)$ et $\mathrm{HOT}_W(S')$,} en termes d'opérations \deleted{\unsure{entre}} \inserted{sur les} catégories sur $S,S'$, \deleted{\unsure{quel}} comme suggéré à l'instant, qu'en se bornant à des ``modèles'' $X$ qui sont dans $\underline{\mathrm{Fib}}\,S$, \deleted{\unsure{tout}} au moins. Et je ne doute pas en effet que ça marche sur $\underline{\mathrm{Fib}\,S}$ pas moins bien que directement sur $\underline{\mathrm{Fibsc}}\,S=$ en termes de $\underline{\mathrm{Hom}}(S^{\unsure{\mathrm o|\circ}},\mathrm{Cat})$.

\centerline{\(\backslash\)}

Pour les axiomes essentiels sur les dérivateurs, je \unsure{renvois|renvoie} à I p. 67-68 et p. 111-112. On a vu que Der 1, \textcircled{\scriptsize Der 2}\textsuperscript{$*$} sont satisfaits par un prédérivateur défini par un couple $(\mathcal M,W\subset\mathrm{Fl}(\mathcal M))$, donc c'est OK ici, \inserted{pour $\mathrm{HOT}_W$}. Il s'ensuit que c'est OK aussi pour $\mathrm{HOT}^{\unsure{\mathrm{lc}|\mathrm{le}}}_W$, \deleted{qui} vu sa caractérisation \inserted{des $\mathrm{HOT}^{\mathrm{lc}}_W(S)$} en \deleted{\unsure{tous}} par sous-catégories des $\mathrm{HOT}_W(S)$. Pour Dér 1, ça marche \inserted{-il} même pour des sommes quelconques de catégories ?, i.e. a-t-on

% slanted
\margin{$*$ \unsure{Cavalier ça} pour \uline{Der} 2 \unsure{car}, je ne sais si c'est \unsure{tj.} vrai même en supposant $W$ \unsure{très} \uline{fortement} saturé \unsure{cf.}, i.e. tel que $\mathrm{hot}_W(f)$ iso $\Rightarrow f\in W$.}

\[
\mathrm{HOT}_W\bigl(\coprod_i S_i\bigr)\xrightarrow{\approx}\prod_i\mathrm{HOT}_W(S_i).
\]

On a en tous cas

\[
\underline{\mathrm{Fib}}\bigl(\overbrace{\coprod_i S_i}^{S}\bigr)
 \xrightarrow{\approx}\prod_i\mathrm{Fib}(S_i)
\]

et le localiseur $W^f_S$ \deleted{\unsure{en}} dans $\underline{\mathrm{Fib}}(S)$ s'identifie
\scanpage{186}
\mspage{131}
% archive sheet number: 185

bel et bien au produit des localiseurs $W^f_{S_i}$ dans les $\underline{\mathrm{Fib}}\,S_i$. Mais dans le cas d'une famille \inserted{peut-être} infinie de catégories $\unsure{\mathcal F|F}_i$, \unsure{munies|munie} des localiseurs $W_i$, si on pose

\[
\deleted{\unsure{\mathcal F}}\mathcal F=\prod\mathcal F_i,\qquad W=\prod W_i\subset\mathrm{Fl}(\mathcal F),
\]

le foncteur canonique

\[
\mathcal F W^{-1}\unsure{\mathord{-}|\longrightarrow|\longleftarrow}\prod_i(\mathcal F_i W_i^{-1})
\]

est-il toujours une équivalence? \deleted{Même dans} le cas $I$ fini il faut \inserted{(il me semble)} un argument assez soigneux, soit par la propriété universelle par réc. sur card. $I$ (on est ramené au cas de deux facteurs), soit à coups de \deleted{\unsure{catégories}} \inserted{catégories} de chemins, en notant que

\[
\underline{\mathrm{Ch}}^{W}_{\infty}(\inserted{\mathcal F;}(s_i),(t_i))
 \xrightarrow{\sim}\prod_i\underline{\mathrm{Ch}}^{W_i}_{\infty}(\mathcal F_i;s_i,t_i),
\tag*{($\alpha$)}
\]

et en passant aux $\pi_0$, qui donne

\[
\mathrm{Hom}_W((s_i),(t_i))\xrightarrow{\sim}\prod_i\mathrm{Hom}_{W_i}(s_i,t_i).
\]

Mais la formule $\alpha$, visiblement, ne marche que pour $I$ fini, Dans le cas général, on trouve la sous-catégorie \inserted{$\prod_i\unsure{!|{}^\prime}$} de $\prod_i$ formée des systèmes $(\unsure{c_i|C_i})_{i\in I}$ de chemins qui sont de longueur bornée.

Donc la question de la validité de Der 1 sans hyp de finitude sur l'ens. d'indices $I$
\scanpage{187}
\mspage{132}
% archive sheet number: 186

reste en suspens\unsure{.|,} pour l'\unsure{instant}. (Je soupçonne que c'est vrai, mais il faudrait, pour le prouver, tenir une description des flèches des $\mathrm{HOT}_W(S)$ par des chemins de longueur bornée, disons \deleted{\unsure{par}} de longueur deux ou trois. C'est une question intéressante, mais que je réserve pour plus tard, si le besoin m'y contraint\ldots

Plus importante est la question de la validité de Der 3 et de ses variantes. (I p. 111-112) Le minimum qu'on demande d'un dérivateur $\mathbb D$, c'est que si $S'\longrightarrow S$ dans $\mathrm{Diag}$ est une Hot-équivalence, ce soit une $\mathbb D$-équivalence, i.e. \deleted{que} (dans le cas actuel) que le foncteur

\[
f^*:\mathrm{HOT}_W(S)\longrightarrow\mathrm{HOT}_W(S')
\]

soit ple fid. sur la sous-catégorie des objets \uline{constants} de $\mathrm{HOT}_W(S)$. Cette condition devient plus tangible quand on dispose déjà \inserted{soit} des foncteurs cohomologiques $H^{\unsure{\bullet|\circ}}_{\mathbb D}(S,-)$, $H^{\unsure{\cdot|\prime}}_{\mathbb D}(S',-)$, soit de l'homologie $H^{\mathbb D}_{\cdot}(S,-),H^{\mathbb D}_{\cdot}(S',-)$. En termes de cette dernière, elle signifie que

\[
\deleted{\unsure{f}}\quad p_!p^*(a)\unsure{\overset{\sim}{\text{---}}|\xrightarrow{\sim}}p'_!p'^*(a)\qquad\text{iso},
\]
\scanpage{188}
\mspage{133}
% archive sheet number: 187

pour \deleted{tout} $a$ dans $\mathcal A=\mathbb D(e)$. (\deleted{\unsure{D'où}} $p:S\longrightarrow e$, $p':S'\longrightarrow e$ sont les morph. structuraux sur la cat. ponctuelle.) On a vu la formule heuristique générale

\[
p_!p^*\underbrace{(a)}_{a_S}\unsure{(\!=\,\prime|=}\,\mathrm{hot}(S)\times a,
\]

où $\mathrm{hot}\,S$ est la classe de $S\in\mathrm{Ob}\,\mathrm{Cat}$ dans $\mathrm{Hot}^{\unsure{?|S}}=(\mathrm{Cat})W_\infty^{-1}$, et où le produit $\times$ est une opération externe de $\mathrm{Hot}$ sur $\mathcal A=\mathbb D(e)$

\[
\mathrm{Hot}\times\mathcal A\longrightarrow\mathcal A.
\]

Dans le cas actuel, $\mathcal A=\mathrm{Hot}$, et l'opération

\[
\mathrm{Hot}\times\mathrm{Hot}\longrightarrow\mathrm{Hot}
\]

ne doit être autre que le \deleted{produit} foncteur ``produit cartésien'', \inserted{lequel est} déduit du foncteur similaire dans $\mathrm{Cat}$. Quand on dispose de ce formalisme, le résultat cherché pour un $f:S'\longrightarrow S$ se réduit simplement à dire que si $f$ est dans $W$, alors il en est de $=$ de $f\times\mathrm{id}_X$ pour tout $X$ dans $\mathrm{Cat}$ (en prenant $X$ tel que $\mathrm{hot}(X)\simeq a$). Or là c'est un fait bien connu, résultant des axiomes mis sur $W$.

J'ai donc l'impression que je n'ai aucune chance de prouver Der 3 a$'$ (\unsure{-p.|p.} 111), ni même Der 3 a) (cas où $S=e$), sans avoir une maîtrise du formalisme des $p_!$. Et j'ai
\scanpage{189}
\mspage{134}
% archive sheet number: 188

bon espoir que cette maîtrise va s'obtenir à bon compte, alors que la question des $p_*$ semble beaucoup plus délicate.

Je vais donc me tourner d'abord vers l'axiome Der 4$'$ d'existence des $\deleted{\unsure{f_!}}$

\[
f_!:\mathrm{HOT}_W(S')\longrightarrow\mathrm{HOT}_W(S),
\]

comme adjoint à gauche de $f^*$. Je vais essayer de construire ad-hoc un foncteur $f_!$, puis \deleted{\unsure{et les}} les deux flèches d'adjonction

\[
\begin{aligned}
\mathrm{id}&\longrightarrow f^* f_!\\
\mathrm{id}&\longleftarrow f_! f^*,
\end{aligned}
\]

\unsure{et} enfin, de vérifier les deux compatibilités pertinentes.

L'idée géométrique ``évidente'', c'est pour $X'\in\mathrm{Cat}_{/S'}$, \deleted{\unsure{dé}}représentant $\xi'\in\mathrm{Ob}\,\mathrm{HOT}_W(S')$, de définir $f_!(\xi')$ par

\[
f_!(\mathrm{hot}_{S',W^{\unsure{*|\prime}}}(\deleted{\unsure{\xi}}X'))=\mathrm{hot}_{S,W}(X')
\]

où dans le second membre, $X'$ est regardé comme objet de $\mathrm{Cat}_{/S}$, grâce à $f:S'\longrightarrow S$ et au foncteur

\[
f_!:\mathrm{Cat}_{/S'}\longrightarrow\mathrm{Cat}_{/S}
\]

qu'il induit. Le seul problème, c'est de voir si ce foncteur est compatible avec les
\scanpage{190}
\mspage{135}
% archive sheet number: 189

localiseurs $W^d_{S'},W^d_S$, et sinon, s'il l'est tout au moins sur la sous-catégorie $\underline{\mathrm{Fib}}\,S'$ ou $\underline{\mathrm{Fib}}\,\underline{\mathrm{sc}}\,S'$, qui suffit pour calculer $\mathrm{HOT}_W(S')$ par localisation. Il faut donc voir si un $X'_1\longrightarrow X'_{\unsure{2|3}}$ dans $\mathrm{Cat}_{/S'}$, qui est dans $W_{S'}^d$ i.e. qui est dans $W$ et le reste par tt chgt de base \deleted{qui est} de la forme $ {}_{s'}\!\backslash S'\longrightarrow S'$, plus gén. par tt \inserted{co-}isom. local, est aussi dans $W_S^d$, i.e. si les chgts de base de la forme $ {}_{s}\!\unsure{\backslash|\prime} S\longrightarrow S$ le transforment en $W$-équivalences (grossières). Or les \deleted{chgt de base} coisom. locaux sont stables par \uline{tt chgt de base} (cf I, p. 55 dans le cas \inserted{dual} des isom locaux), donc on gagne !

\diagram{\begin{tikzcd}[ampersand replacement=\&]
\overline{S'}={}_{s'}\!\backslash S'\arrow[r]\arrow[d]\&S'\arrow[d]\\
\overline S_{\deleted{\unsure{s}}}={}_{s}\!\backslash S\arrow[r]\&S
\end{tikzcd}}

Je vais quand même le \unsure{vérifier|réexpliquer} encore. Si

\diagram{\begin{tikzcd}[ampersand replacement=\&]
\overline{S'}\arrow[r]\arrow[d]\&\unsure{s'|S'}\arrow[d,dash]\\
\overline S\arrow[r,dash]\&S
\end{tikzcd}}

est un carré cartésien dans $(\mathrm{Cat})$, on a dit dans loc. cit., essentiellement, que pour $\overline{s'}$ dans $\overline{S'}$, ayant des images $s',\bar s,\unsure{\bar s|s}$ dans $S',\overline S,S$ \unsure{respt.}, le carré des colocalisés

\diagram{\begin{tikzcd}[ampersand replacement=\&]
{}_{\overline{s'}}\!\backslash\overline{S'}\arrow[r]\arrow[d]\&{}_{s'}\!\backslash S'\arrow[d]\\
{}_{\bar s}\!\backslash\overline S\arrow[r]\&{}_{s}\!\backslash S
\end{tikzcd}}

est encore
\scanpage{191}
\mspage{136}
% archive sheet number: 190

cartésien. Je viens encore de le revérifier à l'instant, pour les objets tout au moins. \deleted{Si} \inserted{Dire que} $\overline S\longrightarrow S$ est un isom local, signifie que les $ {}_{\bar s}\!\backslash\overline S\longrightarrow {}_{s}\!\backslash S$ sont iso. Mais le carré cartésien précédent montre qu'alors les $ {}_{\overline{s'}}\!\backslash\overline{S'}\unsure{\mathord{-}|\longrightarrow} {}_{s'}\!\backslash S'$ sont aussi iso, on gagne !

On a donc le

\uline{\deleted{Proposition} \inserted{Théorème} 1} Pour tte flèche $f:S'\longrightarrow S$ de $\mathrm{Cat}$, le foncteur induit

\[
f_!:\mathrm{Cat}_{/S'}\longrightarrow\mathrm{Cat}_{/S}
\]

est compatible avec les localiseurs $(W^d_{S'},W^d_S)$, donc dualement avec les localiseurs $(W^g_{S'},W^g_S)$, et trivialement avec les localiseurs $(W_S,W_S)$. Il induit donc un diagramme commutatif de foncteurs sur les localisés

\diagram{\begin{tikzcd}[ampersand replacement=\&,column sep=large,row sep=large]
\mathrm{HOT}^{\unsure{l|\,}}_W(S')\arrow[r,"\rho_{S'}"]\arrow[d,"f_!"] \& \mathrm{HOT}^{\unsure{\mathrm{lc}|\mathrm{le}}}_W(S')\arrow[d,"{f_!^{\mathrm{lc}}}"] \& \mathrm{HOT}^{\unsure{\mathrm o|\circ}}_W(S')\arrow[l,"\sigma_{S'}"']\arrow[d,"{f_!^{\mathrm o}}"']\\
\mathrm{HOT}_W(S)\arrow[r,"\rho_S"] \& \mathrm{HOT}^{\mathrm{lc}}_W(S) \& \mathrm{HOT}^{\mathrm o}_W(S)\arrow[l,"\sigma_S"']
\end{tikzcd}.}

\uline{Déf 1} On a défini ainsi les foncteurs verticaux $f_!,f_!^{\mathrm{lc}},f_!^{\mathrm o}=(f^{\mathrm o})_!$, appelés foncteurs images directes homologiques, ou foncteurs d'intégration pour $f$ (et les types d'homotopie \unsure{relatifs|relts/$S$}).
\scanpage{192}
\mspage{137}
% archive sheet number: 191

On a obtenu en même temps, gratuitement, la compatibilité de ces foncteurs aux foncteurs localisation $\rho_{S'},\sigma_{S'}$ \deleted{sur} pour $\mathrm{HOT}(S')$, et $\rho_S,\sigma_S$ pour $\mathrm{HOT}(S)$. Par contre, on ne s'attendra pas qu'ils commutent avec les foncteurs pl. fid. d'inclusion de $\mathrm{HOT}^{\unsure{\mathrm{lc}|\mathrm{le}}}_W$ dans $\mathrm{HOT}_W$. En effet, si on ne suppose pas \uline{$f$ dans Fibhot},

% slanted
\margin{ce cas doit être regardé avec soin ultérieurement.}

il est tout à fait immoral de s'attendre à ce que $f_!$ de $\mathrm{HOT}_W(S')$ transforme types loc constants en item. La situation est donc exactement inverse de celle pour le foncteur $f^*$, qui était compatible avec les foncteurs d'inclusion $\unsure{i_{S'}|i_S},i_{S'}$ et $\unsure{j_S,j_{S'}|j_{S'},j_{S'}}$ -- c'était à tel point trivial qu'il n'a pas été question même de le dire, voire en faire une proposition ! La situation est même encore meilleure pour le foncteur $f_!$ entre les $\mathrm{HOT}$, puisque la description de ce foncteur peut se faire directement sur $\mathrm{Cat}_{/S'}$ tout entier, sans avoir à se borner à des sous-catégories plus ou moins
\scanpage{193}
\mspage{138}
% archive sheet number: 192

\unsure{astucieuses}, où le foncteur induit serait compatible avec les localiseurs ($\underline{\mathrm{Liss}}\,S$ et $\underline{\mathrm{Fibhot}}\,S$, dans le cas de $f^*$).

On a de façon immédiate la transitivité de $f_!,f_!^{\unsure{\mathrm{lc}|\mathrm{le}}},f_!^{\unsure{\mathrm o|\circ}}$, \unsure{sc.|scil.}

\[
(gf)_!=g_!f_!,\quad (gf)_!^{\mathrm{lc}}=g_!^{\mathrm{lc}}f_!^{\mathrm{lc}},\quad
(gf)_!^{\mathrm o}=g_!^{\mathrm o}f_!^{\mathrm o}.
\]

Mais on ne va pas faire bouger $S,S',f$ pour le moment, car il faut maintenant, avant toute autre chose, essayer d'établir l'adjonction de $f_!$ et de $f^*$ (dans cet ordre).

Pour le moment, occupons-nous de $\mathrm{HOT}_W$ seulement -- le cas de $\mathrm{HOT}_W^{\mathrm o}$ s'en déduira, en appliquant ce qu'on saura à $f^{\mathrm o}$. Et on s'occupera de $\mathrm{HOT}_W^{\mathrm{lc}}$ plus tard.

\uline{1)} L'hom

\[
\alpha:\mathrm{id}_{\mathrm{HOT}_W(S')}\longrightarrow f^*f_!
\]

Si $X'$ est dans $\mathrm{Cat}_{/S'}$, $f_!(\mathrm{hot}_{\unsure{S',W}}(X'))$ est $\mathrm{hot}_{S,W}(X')$, mais on ne peut trouver $\deleted{\unsure{f^*}}$ son image inverse $f^*$ sur $S'$, en prenant l'image inverse catégorique. Car même si on fait des hypothèses draconiennes sur $X'$ en tant qu'objet de $\mathrm{Cat}_{/S'}$ (p. ex. $X'$ fibré sur $S'$), celles-ci seront réduites à néant quand on aura composé avec $f$, sur lequel nous ne voulons ici faire aucune hypothèse. Il
\scanpage{194}
\mspage{139}
% archive sheet number: 193

faut donc d'abord appliquer le foncteur $\unsure{\Phi|\mathscr Q}_S$ à $\unsure{\,/\,}X'$ sur $S$, puis prendre son image inverse sur $S'$ :

\[
f^*f_!(\mathrm{hot}_{S',W}(X'))=\mathrm{hot}_{S',W}\bigl(\Phi_S(X'/S)\times_S S'\bigr).
\]

Il faut donc trouver une flèche

\[
\alpha_{X'}: X'\longrightarrow \Phi_S(X'/S)\times_S S'
\]

dans le localisé $(\mathrm{Cat}_{/S'})(W^d_{S'})^{-1}$, voire même une flèche dans $\mathrm{Cat}_{/S'}$ lui-même. Mais par les propriétés d'adjonction des foncteurs $f_!$ et $f^*$ au \uline{niveau catégorique}, une telle flèche revient à une flèche

\[
f_!(X')\longrightarrow\Phi_S(X'/S)
\]

dans $\mathrm{Cat}_{/S}$, i.e. à une flèche qui rend commutatif le carré

\diagram{\begin{tikzcd}[ampersand replacement=\&]
X'\arrow[r]\arrow[d]\&\Phi_S(X'/\unsure{S)|S'|S})\arrow[d]\\
S'\arrow[r]\&S
\end{tikzcd}.}

Or on a une telle flèche canonique, savoir $i_{X'/S}$. Elle est fonctorielle en $X'$ sur $S$ (sans avoir $=$\unsure{.} à se soucier si $X'$ \deleted{\unsure{pour}}\unsure{sur} $S$ provient d'un $X'$ sur $S'$, i.e. si la flèche structurale $X'\longrightarrow S$ se factorise par $S'\longrightarrow\unsure{S|s}$.)
